%%%%%%%%%%%%%%%%%%%%%%%%%%%%%%%%%%%%%%%%%%%%%%%%%%
% redefinition of command
% history 
% v1.0 : 27 Feb, 2026
%%%%%%%%%%%%%%%%%%%%%%%%%%%%%%%%%%%%%%%%%%%%%%%%%%

% kana counter (using expl3)
\ExplSyntaxOn
\cs_new:Npn \kana_map:n #1 {
  \ifcase #1 \or ア \or イ \or ウ \or エ \or オ \or カ \or キ \or ク \or ケ \or コ
  \or サ \or シ \or ス \or セ \or ソ \or タ \or チ \or ツ \or テ \or ト 
  \or ナ \or ニ \orヌ \or ネ \or ノ \or ハ \or ヒ \or フ \or ヘ \or ホ
  \or マ \or ミ \or ム \or メ \or モ \or ヤ \or ユ \or ヨ \or ラ \or リ
  \or ル \or レ \or ロ \or ワ \or ヲ \or ン \fi
}
\NewDocumentCommand{\kana}{m}{
  \expandafter \kana_map:n \value{#1}
}
\ExplSyntaxOff
\AddEnumerateCounter{\kana}{\kana_map:n}{ア}
%\RenewDocumentCommand{\emph}{m}{\textbf{\textsf{\textgt{#1}}}}
%\RenewDocumentCommand{\emph}{m}{{\sffamily\bfseries\gtfamily\kanjiseries{m}\selectfont #1}}
\DeclareTextFontCommand{\emph}{\sffamily\bfseries\gtfamily\kanjiseries{m}}

%%% mathematics
% change font
\NewDocumentCommand{\mr}{m}{\symup{#1}}
\NewDocumentCommand{\mb}{m}{\symbf{#1}}
\NewDocumentCommand{\bs}{m}{\symbf{#1}}
\NewDocumentCommand{\bb}{m}{\symbb{#1}}
\NewDocumentCommand{\scr}{m}{\symscr{#1}}
%\RenewDocumentCommand{\emph}{m}{\textgt{\textsf{\textbf{#1}}}}

% sets
\NewDocumentCommand{\R}{}{\symbb{R}}
\NewDocumentCommand{\N}{}{\symbb{N}}
\NewDocumentCommand{\C}{}{\symbb{C}}
\NewDocumentCommand{\Z}{}{\symbb{Z}}

% const.
\let\i\relax
\DeclareMathOperator{\i}{\symup{i}}
\DeclareMathOperator{\e}{\symup{e}}
%\RenewDocumentCommand{\i}{}{\symup{i}}
%\NewDocumentCommand{\e}{}{\symup{e}}

% matrix, average, order
\NewDocumentCommand{\mqty}{m}{\begin{matrix} #1 \end{matrix}}
\NewDocumentCommand{\avg}{m}{\langle #1 \rangle}
\NewDocumentCommand{\order}{m}{\mathcal{O}(#1)}
%\NewDocumentCommand{\eval}{m}{
%  #1 \mathclose{
%    \rule[-2.0ex]{0.4pt}{5.5ex}\,
%  }
%}
\NewDocumentCommand{\eval}{ s O{\bigg} m }{%
  \IfBooleanTF{#1}{%
    \left. #3 \right|%
  }{%
    \IfValueTF{#2}{%
      \mathopen{#2.} #3 \mathclose{#2\rvert}%
    }{%
      #3 \mathclose{\rvert}%
    }%
  }%
}

% differential op
\NewDocumentCommand{\grad}{}{\nabla}
\RenewDocumentCommand{\div}{}{\nabla \cdot}
\NewDocumentCommand{\rot}{}{\nabla \times}
\NewDocumentCommand{\nb}{}{\nabla}

%% DIFFCOEFF PKG SETTING
\difdef { l } { dn } { style = d^ }
\NewDocumentCommand \dn { m m } { \dl.dn.[#1]{#2} }
\difdef{l}{p}{op-symbol=\partial}
\NewDocumentCommand \dlp {} { \dl.p. }
%\difdef { l } {} {outer-Rdelim = \,}
\difdef { l } {} { outer-Rdelim = \mathchoice{\,}{}{}{} }
\difdef {f, s} {D}{op-symbol = \symup{D}}
\NewDocumentCommand{\Dl}{m}{\upDelta{#1}}

%%% abbrv.
\NewDocumentCommand{\eps}{}{\varepsilon}

%%% layout
\NewDocumentCommand{\qq}{}{\quad}
\NewDocumentCommand{\qqtext}{m}{\quad \text{#1} \quad}
\NewDocumentCommand{\ds}{}{\displaystyle}

%%% color
\definecolor{mint}{HTML}{3EB370}
\definecolor{dark-mint}{HTML}{2E8555}
\RenewDocumentCommand{\CancelColor}{}{\color{ForestGreen}}
\NewDocumentCommand{\blue}{m}{\textcolor{blue}{#1}}
\NewDocumentCommand{\red}{m}{\textcolor{red}{#1}}

%%% other configuration
%\RenewDocumentCommand{\qedsymbol}{}{\blacksquare}
\AtBeginDocument{
\RenewDocumentCommand{\figurename}{}{Figure\,}
\RenewDocumentCommand{\proofname}{}{\bfseries\sffamily\gtfamily Proof.}
}





% equation number (section.eq)
\numberwithin{equation}{section}


\RenewPageStyle{plain}{ % jlreq標準のplainスタイル（ページ番号のみ）を更新
  nombre=\thepage,
  nombre_position=top-right,
  }

% 定義したスタイルを適用
\pagestyle{plain}
